\chapter{Randomized Data Structures: The Dictionary Problem and Skip Lists}
\label{chap:randomized-data-structures}

\begin{center}
  \textbf{Lecture 9}\\
  \emph{Randomized Data Structures: The Dictionary Problem, Block Decomposition, and Skip Lists with Backward Analysis}\\[0.5em]
  Date: October 12, 2026
\end{center}

In previous lectures, we investigated hierarchical memory models (\textit{External Memory Model}) and high-velocity data streams (\textit{Data Streams}), designing algorithms capable of minimizing I/O transfers and sampling elements under severe memory constraints. Across all these domains, algorithmic efficiency relies critically on maintaining dynamic collections of elements through associative structures that support search, insertion, and deletion queries in sublinear or logarithmic time.

In this chapter, we explore a foundational architectural principle in algorithm engineering: \textbf{randomization applied to data structures}. Instead of enforcing rigid deterministic balance invariants (such as complex tree rotations in AVL trees or re-coloring cases in Red-Black trees), we introduce pseudo-random coin tosses directly inside the algorithms. We demonstrate how randomization decouples asymptotic execution performance from input data distributions, achieving logarithmic guarantees $\mathcal{O}(\log n)$ both in expectation and \textit{with high probability} (w.h.p.), while enabling simpler, lightweight implementations ideal for multi-core concurrent environments.

We trace the evolution from deterministic binary search over linked lists to the \textbf{Skip List} introduced by William Pugh, formalizing its structural height and conducting a rigorous probabilistic search cost derivation via \textbf{Backward Analysis}. In the subsequent chapter (Chapter~\ref{chap:randomized-search-trees-treaps-en}), we extend this randomized paradigm directly to binary search trees by examining the \textbf{Treap} and its geometric applications.

\FloatBarrier
\section{Introduction and Motivation: The Dictionary Problem}
\label{sec:dictionary-problem-and-motivation}

\subsection{Formalization of the Dictionary Problem}
In theoretical computer science and high-performance systems engineering, the \textbf{dictionary problem} models the dynamic maintenance of a finite subset $S$ drawn from a totally ordered universe $\mathcal{U}$, where each item is identified by a unique key $k \in \mathcal{U}$.

\begin{definition}[Dynamic Dictionary Problem]
Let $\mathcal{U}$ be a totally ordered universe equipped with a comparison relation $\le$, and let $S \subseteq \mathcal{U}$ be a dynamic subset of cardinality $|S| = n$. A dictionary data structure maintains $S$ subject to an arbitrary sequence of the following core operations:
\begin{enumerate}
  \item $\textsc{Search}(k)$: determines whether an item with key $k$ exists in $S$. If present, it returns a pointer/reference to the corresponding element; otherwise, it returns a null sentinel value ($\textsc{NIL}$).
  \item $\textsc{Insert}(x)$: adds a new item $x$ with key $k = \operatorname{Key}(x)$ into the set, updating $S \leftarrow S \cup \{x\}$.
  \item $\textsc{Delete}(x)$: removes item $x$ from $S$, updating $S \leftarrow S \setminus \{x\}$.
\end{enumerate}
Beyond these three fundamental operations, real-world systems (such as DBMS storage engines and indexing layers) require efficient support for \textbf{order-related queries}:
\begin{enumerate}
  \setcounter{enumi}{3}
  \item $\textsc{Predecessor}(k)$: returns the element in $S$ with the greatest key strictly smaller than $k$:
  \begin{equation}
    \operatorname{pred}(k) = \operatorname{argmax}_{x \in S} \{\operatorname{Key}(x) \mid \operatorname{Key}(x) < k\}.
  \end{equation}
  \item $\textsc{Successor}(k)$: returns the element in $S$ with the smallest key strictly greater than $k$:
  \begin{equation}
    \operatorname{succ}(k) = \operatorname{argmin}_{x \in S} \{\operatorname{Key}(x) \mid \operatorname{Key}(x) > k\}.
  \end{equation}
\end{enumerate}
\end{definition}

\subsection{Comparison of Traditional Deterministic Data Structures}
Classical algorithm design provides distinct trade-offs to resolve the dictionary problem. Table~\ref{tab:deterministic-dictionary-comparison-en} summarizes the asymptotic worst-case computational complexity of the primary deterministic data structures.

\begin{table}[htbp]
\centering
\footnotesize
\setlength{\tabcolsep}{3.5pt}
\begin{tabularx}{\textwidth}{@{}lccccY@{}}
\toprule
\textbf{Data Structure} & \textbf{Search} & \textbf{Insert} & \textbf{Delete} & \textbf{Pred/Succ} & \textbf{Implementation Notes} \\
\midrule
\textbf{Sorted Array} & $\mathcal{O}(\log n)$ & $\mathcal{O}(n)$ & $\mathcal{O}(n)$ & $\mathcal{O}(\log n)$ & Excellent binary search on static data; mutations require linear shifts of contiguous memory. \\
\addlinespace
\textbf{Sorted Linked List} & $\mathcal{O}(n)$ & $\mathcal{O}(1)^*$ & $\mathcal{O}(1)^*$ & $\mathcal{O}(n)$ & Local $\mathcal{O}(1)$ updates given a pointer; search is strictly sequential. \\
\addlinespace
\textbf{Balanced BSTs} & $\mathcal{O}(\log n)$ & $\mathcal{O}(\log n)$ & $\mathcal{O}(\log n)$ & $\mathcal{O}(\log n)$ & Optimal worst-case bounds, but high structural overhead: color/height bits, multiple rotations, coarse locking. \\
(AVL, Red-Black) & & & & & \\
\bottomrule
\end{tabularx}
\caption{Asymptotic worst-case comparison of deterministic data structures ($^*$local cost given pointer).}
\label{tab:deterministic-dictionary-comparison-en}
\end{table}

\subsection{Why Introduce Randomization?}
The deterministic paradigm exhibits two critical architectural bottlenecks in modern systems:
\begin{enumerate}
  \item \textbf{Restructuring Overhead and Implementation Complexity:} Enforcing strict balance invariants requires topological checks, extra metadata per node (balance factors in AVL trees, color bits in Red-Black trees), and cascading single/double rotations upon mutations.
  \item \textbf{Concurrency Bottlenecks in Multi-Threaded Engines:} In deterministic balanced trees, a single insertion or deletion can trigger a cascade of rotations propagating from leaves all the way up to the tree root. In concurrent multi-core architectures, this demands acquiring exclusive write locks over large subtrees or even the root itself, severely degrading parallel scalability.
\end{enumerate}

Introducing \textbf{randomization} revolutionizes this design space by letting internal randomized decisions (such as unbiased coin flips or independent uniform random priorities) govern structure topology:
\begin{itemize}
  \item \textbf{Input Independence:} Asymptotic performance does not depend on key orderings or adversarial inputs, but solely on internal pseudo-random choices.
  \item \textbf{Logarithmic Guarantees With High Probability:} The $\mathcal{O}(\log n)$ running times are achieved not merely in expectation $\mathbb{E}[T]$, but \textbf{with high probability} (w.h.p.), meaning with failure probability at most $\mathcal{O}(n^{-c})$ for any chosen constant $c \ge 1$.
  \item \textbf{Lightweight and Lock-Free Concurrency:} Randomized structures rely on horizontal and vertical link pointers, making lock-free concurrent designs with atomic Compare-And-Swap (CAS) primitives straightforward.
\end{itemize}

\paragraph{Industrial Applications.}
\begin{itemize}
  \item \textbf{Skip Lists:} Used as the primary in-memory table (\textit{MemTable}) in log-structured merge-tree (LSM-Tree) storage engines, including Google's \textbf{LevelDB} and Meta's \textbf{RocksDB}. They are also used in \textbf{Redis} for \textit{Sorted Sets} (\texttt{ZSET})~\cite{pugh1990}.
  \item \textbf{Treaps:} Deployed in geometric indexing, ring-based dynamic routing protocols, and spatial range search systems~\cite{seidel1996}.
\end{itemize}

\FloatBarrier
\section{From Binary Search on Lists to Skip Lists}
\label{sec:binary-search-to-skip-list}

In 1990, William Pugh tackled the challenge of replicating binary search efficiency directly over linked lists while preserving fast pointer updates~\cite{pugh1990}.

\subsection{The Base Linked List (Level 0)}
Consider a singly linked sorted list delimited by sentinels $-\infty$ and $+\infty$, storing the set $S = \{3, 6, 7, 10, 15, 16\}$:

\begin{figure}[htbp]
\centering
\begin{tikzpicture}[
    scale=0.85,
    every node/.style={transform shape},
    xscale=1.25,
    listnode/.style={rectangle, draw=navy!80, fill=navy!6, minimum width=0.85cm, minimum height=0.55cm, rounded corners=2pt, font=\small\bfseries, align=center},
    pointer/.style={thick, color=navy!80, ->, >=Stealth}
]
    \node[listnode] (s1) at (0, 0)   {$-\infty$};
    \node[listnode] (n1) at (1.2, 0) {$3$};
    \node[listnode] (n2) at (2.4, 0) {$6$};
    \node[listnode] (n3) at (3.6, 0) {$7$};
    \node[listnode] (n4) at (4.8, 0) {$10$};
    \node[listnode] (n5) at (6.0, 0) {$15$};
    \node[listnode] (n6) at (7.2, 0) {$16$};
    \node[listnode] (s2) at (8.4, 0) {$+\infty$};

    \draw[pointer] (s1) -- (n1);
    \draw[pointer] (n1) -- (n2);
    \draw[pointer] (n2) -- (n3);
    \draw[pointer] (n3) -- (n4);
    \draw[pointer] (n4) -- (n5);
    \draw[pointer] (n5) -- (n6);
    \draw[pointer] (n6) -- (s2);
\end{tikzpicture}
\caption{Base sorted singly linked list with $-\infty$ and $+\infty$ sentinels.}
\label{fig:base-linked-list-en}
\end{figure}

In this structure, any search requires $\mathcal{O}(n)$ time because pointers only advance sequentially one step at a time. Direct access to the median element in constant time is impossible.

\subsection{The Deterministic Two-Level Model}
To accelerate search, we introduce an upper \textbf{express lane} (Level $L_1$) that samples items at fixed step intervals from the base list (Level $L_0$):

\begin{figure}[htbp]
\centering
\begin{tikzpicture}[
    scale=0.85,
    every node/.style={transform shape},
    xscale=1.25, yscale=1.2,
    listnode/.style={rectangle, draw=navy!80, fill=navy!6, minimum width=0.85cm, minimum height=0.55cm, rounded corners=2pt, font=\small\bfseries, align=center},
    pointer/.style={thick, color=navy!80, ->, >=Stealth},
    vpointer/.style={thick, color=gray!60, dashed, ->, >=Stealth}
]
    % Level L1
    \node[listnode] (l1_s1) at (0, 1.2)   {$-\infty$};
    \node[listnode] (l1_n3) at (3.6, 1.2) {$7$};
    \node[listnode] (l1_s2) at (8.4, 1.2) {$+\infty$};

    \node[anchor=east, font=\footnotesize\bfseries, color=navy] at (-0.4, 1.2) {Level $L_1$:};

    \draw[pointer] (l1_s1) -- (l1_n3);
    \draw[pointer] (l1_n3) -- (l1_s2);

    % Level L0
    \node[listnode] (l0_s1) at (0, 0)   {$-\infty$};
    \node[listnode] (l0_n1) at (1.2, 0) {$3$};
    \node[listnode] (l0_n2) at (2.4, 0) {$6$};
    \node[listnode] (l0_n3) at (3.6, 0) {$7$};
    \node[listnode] (l0_n4) at (4.8, 0) {$10$};
    \node[listnode] (l0_n5) at (6.0, 0) {$15$};
    \node[listnode] (l0_n6) at (7.2, 0) {$16$};
    \node[listnode] (l0_s2) at (8.4, 0) {$+\infty$};

    \node[anchor=east, font=\footnotesize\bfseries, color=navy] at (-0.4, 0) {Level $L_0$:};

    \draw[pointer] (l0_s1) -- (l0_n1);
    \draw[pointer] (l0_n1) -- (l0_n2);
    \draw[pointer] (l0_n2) -- (l0_n3);
    \draw[pointer] (l0_n3) -- (l0_n4);
    \draw[pointer] (l0_n4) -- (l0_n5);
    \draw[pointer] (l0_n5) -- (l0_n6);
    \draw[pointer] (l0_n6) -- (l0_s2);

    % Vertical pointers perfectly aligned
    \draw[vpointer] (l1_s1) -- (l0_s1);
    \draw[vpointer] (l1_n3) -- (l0_n3);
    \draw[vpointer] (l1_s2) -- (l0_s2);
\end{tikzpicture}
\caption{Deterministic two-level list: express lane $L_1$ skips groups of items.}
\label{fig:deterministic-two-level-list-en}
\end{figure}

Searching for a key $k$ (e.g., $k = 15$) proceeds in two phases:
\begin{enumerate}
  \item Scan horizontally along $L_1$ from $-\infty$ until finding the greatest item $\le k$ (here, node $7$).
  \item Drop down to level $L_0$ at node $7$ and continue linear local scanning ($7 \to 10 \to 15$).
\end{enumerate}

\paragraph{The Deterministic Dilemma.}
While search is expedited, insertions and deletions disrupt the uniform spacing between express pointers. Maintaining rigid equal intervals requires an expensive $\mathcal{O}(n)$ global reorganization, negating any speedup.

\subsection{Square Root Decomposition}
Let us analytically formalize the query cost under a 2-level deterministic block scheme:
\begin{itemize}
  \item $n$: total number of stored elements;
  \item $k$: number of blocks (and express pointers at level $L_1$);
  \item $\frac{n}{k}$: number of elements per block at the base level $L_0$.
\end{itemize}

\begin{figure}[htbp]
\centering
\resizebox{0.95\textwidth}{!}{%
\begin{tikzpicture}[
    xscale=1.25, yscale=1.2,
    blocknode/.style={rectangle, draw=navy!80, fill=navy!5, minimum width=2.4cm, minimum height=0.9cm, rounded corners=3pt, font=\small, align=center},
    expressnode/.style={rectangle, draw=navy!80, fill=navy!15, minimum width=1.6cm, minimum height=0.6cm, rounded corners=2pt, font=\scriptsize\bfseries, align=center},
    pointer/.style={thick, color=navy!80, ->, >=Stealth},
    vpointer/.style={thick, color=navy!65, dashed, ->, >=Stealth}
]
    % Row labels
    \node[anchor=east, font=\scriptsize\bfseries, color=navy] at (0, 1.8) {Level $L_1$ (Express lane):};
    \node[anchor=east, font=\scriptsize\bfseries, color=navy] at (0, 0) {Level $L_0$ (Base list):};

    % Express lane pointers (k pointers)
    \node[expressnode] (e1) at (1.6, 1.8) {Pointer $1$};
    \node[expressnode] (e2) at (4.8, 1.8) {Pointer $2$};
    \node[expressnode] (e3) at (8.0, 1.8) {Pointer $3$};
    \node[font=\bfseries, color=navy!80] (dots) at (9.8, 1.8) {$\cdots$};
    \node[expressnode] (ek) at (11.4, 1.8) {Pointer $k$};

    % Express pointers
    \draw[pointer] (e1) -- node[above=2pt, font=\scriptsize\bfseries, color=navy] {$\mathcal{O}(k)$ jumps} (e2);
    \draw[pointer] (e2) -- (e3);
    \draw[pointer] (e3) -- (dots);
    \draw[pointer] (dots) -- (ek);

    % Base list blocks (size n/k each)
    \node[blocknode] (b1) at (1.6, 0) {\textbf{Block 1}\\ \scriptsize $n/k$ elements};
    \node[blocknode] (b2) at (4.8, 0) {\textbf{Block 2}\\ \scriptsize $n/k$ elements};
    \node[blocknode] (b3) at (8.0, 0) {\textbf{Block 3}\\ \scriptsize $n/k$ elements};
    \node[font=\bfseries, color=navy!80] (bdots) at (9.8, 0) {$\cdots$};
    \node[blocknode] (bk) at (11.4, 0) {\textbf{Block $k$}\\ \scriptsize $n/k$ elements};

    % Sequential connections between blocks
    \draw[pointer, color=navy!40] (b1) -- (b2);
    \draw[pointer, color=navy!40] (b2) -- (b3);
    \draw[pointer, color=navy!40] (b3) -- (bdots);
    \draw[pointer, color=navy!40] (bdots) -- (bk);

    % Perfectly aligned vertical arrows from express pointers to base blocks
    \draw[vpointer] (e1.south) -- (b1.north);
    \draw[vpointer] (e2.south) -- (b2.north);
    \draw[vpointer] (e3.south) -- (b3.north);
    \draw[vpointer] (ek.south) -- (bk.north);

    % Total elements brace
    \draw[decorate, decoration={brace, amplitude=6pt, mirror}, thick, color=navy!80] (0.4, -0.7) -- (12.6, -0.7)
        node[midway, below=8pt, font=\small\bfseries, color=navy] {$n$ total elements ($k$ blocks of size $n/k$, internal scan cost $\mathcal{O}(n/k)$)};
\end{tikzpicture}%
}
\caption{Deterministic square root decomposition into $k$ blocks of size $n/k$.}
\label{fig:square-root-decomposition-schema-en}
\end{figure}

The worst-case search cost $T(k)$ is the sum of steps along the express lane plus steps inside a block:
\begin{equation}
  T(k) = \mathcal{O}(k) + \mathcal{O}\left(\frac{n}{k}\right) = \mathcal{O}\left(k + \frac{n}{k}\right).
\end{equation}

\begin{theorem}[Optimal Block Size]
The cost function $f(k) = k + \frac{n}{k}$ for $k > 0$ has a unique global minimum at:
\begin{equation}
  k^* = \sqrt{n}.
\end{equation}
At this optimum, the worst-case search complexity is:
\begin{equation}
  T(\sqrt{n}) = \mathcal{O}(\sqrt{n}).
\end{equation}
\end{theorem}

\begin{proof}
Differentiating $f(k) = k + \frac{n}{k}$ with respect to $k$:
\begin{equation}
  \frac{df}{dk} = 1 - \frac{n}{k^2}.
\end{equation}
Setting the first derivative to zero yields:
\begin{equation}
  1 - \frac{n}{k^2} = 0 \iff k^2 = n \iff k = \sqrt{n}.
\end{equation}
The second derivative is strictly positive for all $k > 0$:
\begin{equation}
  \frac{d^2f}{dk^2} = \frac{2n}{k^3} > 0.
\end{equation}
Thus, $f(k)$ is strictly convex and achieves its minimum at $k^* = \sqrt{n}$, giving $f(\sqrt{n}) = 2\sqrt{n} = \mathcal{O}(\sqrt{n})$.
\end{proof}

\begin{quote}
  \textbf{The Key Insight of Skip Lists:} Avoid rigid deterministic reorganization by assigning each node's level via independent probabilistic coin tosses!
\end{quote}

\FloatBarrier
\section{Skip Lists: Structure and Formal Probabilistic Analysis}
\label{sec:skip-list-probabilistic-analysis}

\subsection{Formal Definition}
A Skip List is a sequence of sorted linked lists $\{L_0, L_1, \dots, L_h\}$ satisfying:
\begin{definition}[Skip List]
Let $S$ be a set of $n$ distinct keys. A Skip List for $S$ satisfies:
\begin{enumerate}
  \item Level $L_0$ contains all $n$ keys in ascending order, bounded by sentinels $-\infty$ and $+\infty$.
  \item For each level $i \ge 0$, $L_{i+1} \subset L_i$.
  \item \textbf{Independent Geometric Growth:} An element in $L_i$ is promoted to $L_{i+1}$ independently with probability $p = \frac{1}{2}$ (a fair coin flip yielding \textbf{Heads}). If the flip yields \textbf{Tails}, its tower stops.
\end{enumerate}
\end{definition}

\begin{figure}[htbp]
\centering
\begin{tikzpicture}[
    scale=0.85,
    every node/.style={transform shape},
    xscale=1.25, yscale=1.3,
    listnode/.style={rectangle, draw=navy!80, fill=navy!6, minimum width=0.85cm, minimum height=0.55cm, rounded corners=2pt, font=\scriptsize\bfseries, align=center},
    pointer/.style={thick, color=navy!80, ->, >=Stealth},
    vpointer/.style={thick, color=gray!60, dashed, ->, >=Stealth}
]
    % Level L3 (y = 3.0)
    \node[listnode] (l3_inf)  at (0, 3.0)   {$-\infty$};
    \node[listnode] (l3_15)   at (6.0, 3.0) {$15$};
    \node[listnode] (l3_pinf) at (8.4, 3.0) {$+\infty$};
    \node[anchor=east, font=\scriptsize\bfseries, color=navy] at (-0.4, 3.0) {Level $L_3$:};
    \draw[pointer] (l3_inf) -- (l3_15);
    \draw[pointer] (l3_15) -- (l3_pinf);

    % Level L2 (y = 2.0)
    \node[listnode] (l2_inf)  at (0, 2.0)   {$-\infty$};
    \node[listnode] (l2_3)    at (1.2, 2.0) {$3$};
    \node[listnode] (l2_15)   at (6.0, 2.0) {$15$};
    \node[listnode] (l2_pinf) at (8.4, 2.0) {$+\infty$};
    \node[anchor=east, font=\scriptsize\bfseries, color=navy] at (-0.4, 2.0) {Level $L_2$:};
    \draw[pointer] (l2_inf) -- (l2_3);
    \draw[pointer] (l2_3) -- (l2_15);
    \draw[pointer] (l2_15) -- (l2_pinf);

    % Level L1 (y = 1.0)
    \node[listnode] (l1_inf)  at (0, 1.0)   {$-\infty$};
    \node[listnode] (l1_3)    at (1.2, 1.0) {$3$};
    \node[listnode] (l1_10)   at (4.8, 1.0) {$10$};
    \node[listnode] (l1_15)   at (6.0, 1.0) {$15$};
    \node[listnode] (l1_pinf) at (8.4, 1.0) {$+\infty$};
    \node[anchor=east, font=\scriptsize\bfseries, color=navy] at (-0.4, 1.0) {Level $L_1$:};
    \draw[pointer] (l1_inf) -- (l1_3);
    \draw[pointer] (l1_3) -- (l1_10);
    \draw[pointer] (l1_10) -- (l1_15);
    \draw[pointer] (l1_15) -- (l1_pinf);

    % Level L0 (y = 0.0)
    \node[listnode] (l0_inf)  at (0, 0.0)   {$-\infty$};
    \node[listnode] (l0_3)    at (1.2, 0.0) {$3$};
    \node[listnode] (l0_6)    at (2.4, 0.0) {$6$};
    \node[listnode] (l0_7)    at (3.6, 0.0) {$7$};
    \node[listnode] (l0_10)   at (4.8, 0.0) {$10$};
    \node[listnode] (l0_15)   at (6.0, 0.0) {$15$};
    \node[listnode] (l0_16)   at (7.2, 0.0) {$16$};
    \node[listnode] (l0_pinf) at (8.4, 0.0) {$+\infty$};
    \node[anchor=east, font=\scriptsize\bfseries, color=navy] at (-0.4, 0.0) {Level $L_0$:};
    \draw[pointer] (l0_inf) -- (l0_3);
    \draw[pointer] (l0_3) -- (l0_6);
    \draw[pointer] (l0_6) -- (l0_7);
    \draw[pointer] (l0_7) -- (l0_10);
    \draw[pointer] (l0_10) -- (l0_15);
    \draw[pointer] (l0_15) -- (l0_16);
    \draw[pointer] (l0_16) -- (l0_pinf);

    % Vertical pointers perfectly aligned between .south and .north
    \draw[vpointer] (l3_inf.south)  -- (l2_inf.north);
    \draw[vpointer] (l2_inf.south)  -- (l1_inf.north);
    \draw[vpointer] (l1_inf.south)  -- (l0_inf.north);

    \draw[vpointer] (l2_3.south)    -- (l1_3.north);
    \draw[vpointer] (l1_3.south)    -- (l0_3.north);

    \draw[vpointer] (l1_10.south)   -- (l0_10.north);

    \draw[vpointer] (l3_15.south)   -- (l2_15.north);
    \draw[vpointer] (l2_15.south)   -- (l1_15.north);
    \draw[vpointer] (l1_15.south)   -- (l0_15.north);

    \draw[vpointer] (l3_pinf.south) -- (l2_pinf.north);
    \draw[vpointer] (l2_pinf.south) -- (l1_pinf.north);
    \draw[vpointer] (l1_pinf.south) -- (l0_pinf.north);
\end{tikzpicture}
\caption{Multi-level Skip List representation across 4 levels ($L_0$ through $L_3$).}
\label{fig:multilevel-skip-list-en}
\end{figure}

\subsection{Core Algorithms}

\subsubsection{1. Search ($\textsc{Search}(k)$)}
The search algorithm (Algorithm~\ref{alg:skip-list-search-en}) starts at the highest level $h$ at $-\infty$. If the forward key is $\le k$, it moves right; otherwise, it drops down one level.

\begin{algorithm}[htbp]
\caption{Skip List Search ($\textsc{Search}(k)$)}
\label{alg:skip-list-search-en}
\KwInput{Skip List $SL$ with maximum height $h$, search key $k \in \mathcal{U}$}
\KwOutput{Node containing $k$, or $\textsc{NIL}$ if not present}

$curr \leftarrow \text{head of Skip List at level } h$\;
\For{$i \leftarrow h$ \KwTo $0$}{
    \While{$curr.\text{next}[i] \ne \textsc{NIL}$ \textbf{and} $curr.\text{next}[i].\text{key} \le k$}{
        $curr \leftarrow curr.\text{next}[i]$\;
    }
}
\If{$curr.\text{key} == k$}{
    \Return $curr$\;
}
\Return $\textsc{NIL}$\;
\end{algorithm}

\subsubsection{2. Insertion ($\textsc{Insert}(x)$)}
To insert key $k$:
\begin{enumerate}
  \item Search for the key's position, storing predecessors in array $\text{update}[0 \dots h]$.
  \item Sample a tower height $l$ by simulating independent coin flips:
  \begin{equation*}
    l \leftarrow 0; \quad \textbf{while } \operatorname{random}() < 0.5 \textbf{ do } l \leftarrow l + 1.
  \end{equation*}
  \item Allocate the new node of height $l$ and splice pointers for each level $i \in [0, l]$:
  \begin{equation}
    \text{node}.\text{next}[i] \leftarrow \text{update}[i].\text{next}[i], \qquad \text{update}[i].\text{next}[i] \leftarrow \text{node}.
  \end{equation}
\end{enumerate}

\subsubsection{3. Deletion ($\textsc{Delete}(x)$)}
Find predecessors across all levels and bypass: $\text{update}[i].\text{next}[i] \leftarrow x.\text{next}[i]$, then deallocate.

\subsection{Operational Costs}
\begin{align}
  \operatorname{Cost}(\textsc{Search}) &= \# \text{vertical pointers} + \# \text{horizontal pointers}, \\
  \operatorname{Cost}(\textsc{Insert}) &= \operatorname{Cost}(\textsc{Search}) + \operatorname{height}(\text{key}), \\
  \operatorname{Cost}(\textsc{Delete}) &= \operatorname{Cost}(\textsc{Search}) + \# \text{pointer updates}.
\end{align}

\subsection{Proof: Maximum Height of a Skip List}
\begin{theorem}[Logarithmic Height Bound]
Let $S$ be a set of $n$ keys in a Skip List with $p = \frac{1}{2}$, and let $H = \max_{k \in S} \{\operatorname{level}(k)\}$. For any constant $c > 1$:
\begin{equation}
  \mathbb{P}(H \ge c \log_2 n) \le \frac{1}{n^{c-1}} = n^{-(c-1)}.
\end{equation}
Thus, Skip List height is $\mathcal{O}(\log n)$ with high probability (w.h.p.).
\end{theorem}

\begin{proof}
The proof proceeds in three analytical steps:

\textbf{Step 1: Single element tower height.}
Let $\operatorname{level}(k)$ be the level assigned to an item with key $k$. An element reaches a level $\ge l$ if and only if it obtains at least $l$ consecutive successes (Heads) in independent coin tosses:
\begin{equation}
  \mathbb{P}(\operatorname{level}(k) \ge l) = \left(\frac{1}{2}\right)^l.
\end{equation}

\textbf{Step 2: Union bound across all $n$ keys.}
The overall Skip List height $H$ satisfies $H \ge l$ if and only if at least one key reaches level $\ge l$:
\begin{equation}
  \mathbb{P}(H \ge l) = \mathbb{P}\left(\bigcup_{k \in S} \{\operatorname{level}(k) \ge l\}\right) \le \sum_{k=1}^n \mathbb{P}(\operatorname{level}(k) \ge l) = n \left(\frac{1}{2}\right)^l = \frac{n}{2^l}.
\end{equation}

\textbf{Step 3: Setting logarithmic height threshold $l = c \log_2 n$.}
For any constant $c > 1$:
\begin{equation}
  \mathbb{P}(H \ge c \log_2 n) \le \frac{n}{2^{c \log_2 n}} = \frac{n}{\left(2^{\log_2 n}\right)^c} = \frac{n}{n^c} = \frac{1}{n^{c-1}} = n^{-(c-1)}.
\end{equation}
\end{proof}

\subsection{Search Cost: Backward Analysis}
To determine horizontal traversal steps, we trace the search path \textbf{in reverse}, starting from the target at $L_0$ up to the head sentinel.

\begin{figure}[htbp]
\centering
\resizebox{0.82\textwidth}{!}{%
\begin{tikzpicture}[
    xscale=1.3, yscale=1.3,
    nodebox/.style={rectangle, draw=navy!80, fill=navy!8, minimum width=1.3cm, minimum height=0.7cm, rounded corners=3pt, font=\bfseries, align=center},
    actionarrow/.style={very thick, color=red!75!black, ->, >=Stealth},
    guideline/.style={dashed, color=gray!35, thin}
]
    % Level guidelines
    \draw[guideline] (-5.0, 2.4) -- (2.6, 2.4) node[right, font=\scriptsize\bfseries, color=navy!70] {Level $L_{i+1}$};
    \draw[guideline] (-5.0, 0)   -- (2.6, 0)   node[right, font=\scriptsize\bfseries, color=navy!70] {Level $L_i$};

    % Nodes
    \node[nodebox] (top)  at (0, 2.4)  {$x$};
    \node[nodebox] (left) at (-4.0, 0) {$x'$};
    \node[nodebox] (curr) at (0, 0)    {$x$};

    % Vertical transition (Heads)
    \draw[actionarrow] (curr.north) -- (top.south)
        node[midway, right=14pt, align=left, font=\scriptsize] {
            \textbf{Heads} $\left(p = \frac{1}{2}\right)$\\[3pt]
            Ascend vertically to $L_{i+1}$
        };

    % Horizontal transition to the left (Tails)
    \draw[actionarrow] (curr.west) -- (left.east)
        node[midway, above=10pt, align=center, font=\scriptsize] {
            \textbf{Tails} $\left(1-p = \frac{1}{2}\right)$\\[3pt]
            Step left on $L_i$
        };
\end{tikzpicture}%
}
\caption{Transitions in backward analysis: vertical ascent (Heads) vs horizontal left step (Tails).}
\label{fig:backward-analysis-transitions-en}
\end{figure}

At each step backward from node $x$ at level $i$:
\begin{enumerate}
  \item If $x$ was promoted to $L_{i+1}$ (Heads, $p = 1/2$), the search path entered from $L_{i+1}$; we step vertically up.
  \item If $x$ was not promoted (Tails, $1 - p = 1/2$), the search path came from the left on $L_i$; we step left.
\end{enumerate}

\begin{theorem}[Expected Horizontal Steps per Level]
The expected number of leftward steps before ascending is:
\begin{equation}
  \mathbb{E}[\# \text{tails}] = \sum_{h=0}^{\infty} h \left(\frac{1}{2}\right)^h = 2.
\end{equation}
\end{theorem}

\begin{proof}
For $|x| < 1$, the geometric series evaluates to:
\begin{equation}
  \sum_{h=0}^{\infty} x^h = \frac{1}{1 - x}.
\end{equation}
Differentiating term by term with respect to $x$:
\begin{equation}
  \sum_{h=1}^{\infty} h x^{h-1} = \frac{1}{(1 - x)^2}.
\end{equation}
Multiplying by $x$:
\begin{equation}
  \sum_{h=0}^{\infty} h x^h = \frac{x}{(1 - x)^2}.
\end{equation}
Evaluating at $x = \frac{1}{2}$:
\begin{equation}
  \sum_{h=0}^{\infty} h \left(\frac{1}{2}\right)^h = \frac{\frac{1}{2}}{\left(1 - \frac{1}{2}\right)^2} = \frac{1/2}{1/4} = 2.
\end{equation}
\end{proof}

Since the maximum height is $\mathcal{O}(\log n)$, the total expected search cost is:
\begin{equation}
  \mathbb{E}[\text{Search Cost}] \le 2 \cdot \mathcal{O}(\log n) + \mathcal{O}(\log n) = \mathcal{O}(\log n).
\end{equation}

\subsection{Advanced Application: Vector Databases and HNSW Graphs}
Modern AI vector search applies the Skip List principle to multi-dimensional graphs via \textbf{HNSW} (\textit{Hierarchical Navigable Small World})~\cite{malkov2020}. Dense vectors reside at $L_0$, while upper layers contain sparser subgraphs with long-range links acting as express lanes for logarithmic Approximate Nearest Neighbor Search.

\FloatBarrier
\section{Summary and Outlook: Toward Randomized Search Trees}
\label{sec:skip-list-conclusions-en}

In this chapter we explored how internal algorithmic randomization eliminates the structural rigidity and maintenance overhead of classical balanced search structures:
\begin{itemize}
  \item Starting from the sequential bottleneck of standard linked lists, we derived hierarchical express lanes (transitioning from deterministic $\sqrt{n}$ block partitioning to geometric towers governed by coin-flip promotion probability $p = 1/2$);
  \item The \textbf{Skip List} emulates binary search probabilistically, guaranteeing $\mathcal{O}(n)$ expected space and $\mathcal{O}(\log n)$ runtime with high probability for search, insertion, and deletion;
  \item The \textbf{Backward Analysis} technique inverts the analytical perspective by traversing the query path in reverse from the target key back to the top sentinels, establishing that the expected number of horizontal shifts per level is strictly bounded by $1/p = 2$.
\end{itemize}

While Skip Lists solve the dictionary problem through multi-level linked sequences, an important architectural question arises: \textit{can the same randomized elegance be applied directly to binary search trees, eliminating rigid deterministic balance invariants and complex tree rotations?}
Chapter~\ref{chap:randomized-search-trees-treaps-en} addresses this question by presenting the \textbf{Treap} (Randomized Search Tree), exploring its structural uniqueness, rotation mechanics, the Seidel-Aragon probabilistic analysis, and multi-dimensional geometric applications.
